% ---------------------------------------------------------------------------
% Preâmbulo do relatório — Trabalho 1 de Sistemas Operacionais
% ---------------------------------------------------------------------------

\usepackage{polyglossia}
\setmainlanguage{portuguese}

% ---------------------------------------------------------------- tipografia
\usepackage{fontspec}
\setmainfont{TeX Gyre Pagella}
\setsansfont{TeX Gyre Heros}[Scale=0.92]
\setmonofont{DejaVu Sans Mono}[Scale=0.80]
\usepackage{microtype}

\linespread{1.10}
\widowpenalty=10000
\clubpenalty=10000
\displaywidowpenalty=10000
\setlength{\parskip}{0.55em}
\setlength{\parindent}{0pt}
\setlength{\emergencystretch}{3em}

% -------------------------------------------------------------------- cores
\usepackage{xcolor}
\definecolor{ink}{HTML}{111111}
\definecolor{ink2}{HTML}{3F3E3B}
\definecolor{muted}{HTML}{6E6C66}
\definecolor{rule}{HTML}{D8D7CF}
\definecolor{accent}{HTML}{1C5CAB}
\definecolor{codebg}{HTML}{F6F6F3}
\definecolor{notebg}{HTML}{F2F6FC}
\color{ink}

% ------------------------------------------------------------------ títulos
\usepackage{titlesec}
\titleformat{\section}
  {\sffamily\bfseries\Large\color{ink}}{\color{accent}\thesection}{0.7em}{}
  [\vspace{-0.55em}{\color{rule}\rule{\linewidth}{0.5pt}}]
\titleformat{\subsection}
  {\sffamily\bfseries\large\color{ink}}{\color{accent}\thesubsection}{0.6em}{}
\titleformat{\subsubsection}
  {\sffamily\bfseries\normalsize\color{ink2}}{\color{accent}\thesubsubsection}{0.5em}{}
\titlespacing*{\section}{0pt}{2.0em}{0.8em}
\titlespacing*{\subsection}{0pt}{1.5em}{0.5em}
\titlespacing*{\subsubsection}{0pt}{1.1em}{0.35em}

% ---------------------------------------------------------- cabeçalho/rodapé
\usepackage{fancyhdr}
\usepackage{lastpage}
\pagestyle{fancy}
\fancyhf{}
\renewcommand{\headrulewidth}{0.4pt}
\renewcommand{\headrule}{\hbox to\headwidth{\color{rule}\leaders\hrule height \headrulewidth\hfill}}
\fancyhead[L]{\sffamily\footnotesize\color{muted}Produtor--consumidor com semáforos}
\fancyhead[R]{\sffamily\footnotesize\color{muted}Sistemas Operacionais}
\fancyfoot[C]{\sffamily\footnotesize\color{muted}\thepage\ / \pageref{LastPage}}
\fancypagestyle{plain}{\fancyhf{}\renewcommand{\headrulewidth}{0pt}%
  \fancyfoot[C]{\sffamily\footnotesize\color{muted}\thepage\ / \pageref{LastPage}}}

% ------------------------------------------------------------------ tabelas
\usepackage{booktabs}
\usepackage{array}
\renewcommand{\arraystretch}{1.20}
\setlength{\tabcolsep}{7pt}
\usepackage{caption}
\captionsetup{font={sf,small,color=ink2},labelfont={bf,color=accent},
              skip=6pt,justification=raggedright,singlelinecheck=false}
\captionsetup[table]{position=above,skip=5pt}

% ------------------------------------------------------------------ figuras
\usepackage{graphicx}
\usepackage{float}
\floatplacement{figure}{H}
\usepackage{placeins}
\let\oldsection\section
\renewcommand{\section}{\FloatBarrier\oldsection}

% Um título nunca fica sozinho no pé da página.
\usepackage{needspace}
\usepackage{etoolbox}
\pretocmd{\oldsection}{\Needspace*{6\baselineskip}}{}{}
\pretocmd{\subsection}{\Needspace*{5\baselineskip}}{}{}
\pretocmd{\subsubsection}{\Needspace*{4\baselineskip}}{}{}

% --------------------------------------------------------------------- listas
\usepackage{enumitem}
\setlist[itemize]{leftmargin=1.35em,topsep=0.35em,itemsep=0.30em,parsep=0pt}
\setlist[enumerate]{leftmargin=1.6em,topsep=0.35em,itemsep=0.30em,parsep=0pt}

% ----------------------------------------------------------- caixas de destaque
\usepackage[most]{tcolorbox}
\newtcolorbox{destaque}[1][]{
  enhanced, breakable, sharp corners,
  colback=notebg, colframe=accent,
  boxrule=0pt, leftrule=2.2pt,
  left=11pt, right=11pt, top=9pt, bottom=9pt,
  before skip=1.0em, after skip=1.0em, #1}

% Blockquotes do Markdown viram caixas de destaque.
\renewenvironment{quote}
  {\begin{destaque}\setlength{\parskip}{0.45em}}
  {\end{destaque}}

% ------------------------------------------------------------------- código
\usepackage{fvextra}
\usepackage{framed}
\definecolor{shadecolor}{HTML}{F6F6F3}
\DefineVerbatimEnvironment{Highlighting}{Verbatim}
  {commandchars=\\\{\},fontsize=\footnotesize,breaklines,breakanywhere,xleftmargin=6pt}
\renewenvironment{Shaded}
  {\vspace{0.15em}\begin{snugshade}\setlength{\parskip}{0pt}}
  {\end{snugshade}\vspace{0.15em}}

% --------------------------------------------------------------------- links
\usepackage[hidelinks]{hyperref}
\hypersetup{colorlinks=true,linkcolor=accent,urlcolor=accent,citecolor=accent,
            pdftitle={Produtor-consumidor com semaforos},
            pdfauthor={Pietro Prauchner}}

% --------------------------------------------------------------------- sumário
\usepackage{tocloft}
\renewcommand{\cfttoctitlefont}{\sffamily\bfseries\Large\color{ink}}
\renewcommand{\cftsecfont}{\sffamily\bfseries\color{ink}}
\renewcommand{\cftsubsecfont}{\sffamily\color{ink2}}
\renewcommand{\cftsecpagefont}{\sffamily\bfseries\color{muted}}
\renewcommand{\cftsubsecpagefont}{\sffamily\color{muted}}
\setlength{\cftbeforesecskip}{0.45em}
\renewcommand{\cftdot}{}

% ------------------------------------------------------------------- capa
\newcommand{\capa}{%
\begin{titlepage}
\thispagestyle{empty}
\raggedright
\vspace*{1.2cm}

{\sffamily\footnotesize\color{muted}%
 ENGENHARIA DE SOFTWARE \textbullet\ 4\textordmasculine{} SEMESTRE\par}
\vspace{0.35cm}
{\sffamily\large\color{ink2}Sistemas Operacionais \ --- \ Trabalho 1\par}

\vspace{1.1cm}
{\color{accent}\rule{\linewidth}{2.2pt}}
\vspace{0.9cm}

{\sffamily\bfseries\fontsize{29}{34}\selectfont\color{ink}%
 Exclusão mútua com semáforos\par}
\vspace{0.35cm}
{\sffamily\fontsize{17}{21}\selectfont\color{ink2}%
 Um experimento controlado sobre o problema\\produtor--consumidor com buffer limitado\par}

\vspace{0.9cm}
{\color{rule}\rule{\linewidth}{0.6pt}}

\vspace{1.3cm}
{\itshape\large\color{ink2}%
 Demonstração empírica de que a exclusão mútua \emph{não} é garantida sem
 sincronização explícita, e passa a ser garantida com semáforos POSIX --- com
 medição do custo em tempo de execução.\par}

\vfill

{\sffamily\footnotesize\color{muted}AUTOR\par}
\vspace{0.15cm}
{\sffamily\large\color{ink}Pietro Prauchner\par}

\vspace{0.7cm}
{\color{rule}\rule{\linewidth}{0.6pt}}
\vspace{0.4cm}

{\sffamily\footnotesize\color{muted}%
\begin{tabular}{@{}ll@{}}
Linguagem     & C11 \textbullet\ POSIX \texttt{pthreads} \textbullet\ semáforos \texttt{sem\_t}\\
Ambiente      & WSL2 Ubuntu \textbullet\ gcc 15.2.0 \textbullet\ 6 núcleos\\
Código-fonte  & \texttt{PPrauchner/SO-Trabalho-1-Semaforos}\\
Data          & Setembro de 2026\\
\end{tabular}\par}

\vspace{0.6cm}
\end{titlepage}}
